\paragraph{Reducing the kernel before integration.}
Put $X=x=s-1$ in this paragraph. An integration by parts removes the
poles at $w=\pm\ii$ before any polylogarithmic primitive is constructed.
It combines the elementary and root-dependent parts of the original
soft kernel, so no integrated cross section is needed.

Keep $X=s-1$, $\beta=X/(X+2)$ and $E^2=(X+2)^2/(4s)$, and define
\begin{equation}
 k^2=\frac{1+XU^2}{s},\qquad d=1-\beta^2U^2,\qquad
 a=\operatorname{artanh}(\beta U),\qquad
 f=\operatorname{artanh}(U/k).
\end{equation}
In this section $d$ denotes this rational function, rather than the
spacetime dimension.  On $0<U<1$, all square roots are positive and
both functions are real.  Separate the elementary part of $H$ as
\begin{align}
 h_0&=\overline h_0+\frac{2(X+2)U}{Xk^2},\\
 \overline h_0&=U\left[\frac{X^2+2X+4}{2s}
 -\frac{X^3}{4s^2}-\frac{3X^3(1-U^2)}{8s^2}\right],\\
 h_k&=\frac{b_1}{k}+\frac{2(X+2)(1-U^2)}{Xsk^3},\qquad
 b_1=\frac{X^3-2X^2-16X-8}{X^2}.
\end{align}
The identities
\begin{equation}
 \frac{\dd f}{\dd U}=\frac1{k(1-U^2)},\qquad
 p(U)=\frac{2(X+2)U(1-U^2)}X
\end{equation}
give the exact integration-by-parts relation
\begin{equation}
 \begin{aligned}
 \frac a d\left[\frac{2(X+2)U}{Xk^2}+h_kf\right]
 ={}&\frac{\dd}{\dd U}\left(\frac{apf}{dk}\right)\\
 &+\frac f k\left[
 a\left(\frac{b_1}d+\frac{16sU^2}{X(X+2)d^2}\right)
 -\frac{2U(1-U^2)}{d^2}\right].
 \end{aligned}
 \label{eq:comp-ck-pole-removal}
\end{equation}
Its boundary term vanishes at both ends.  At $U=0$, $a$ and $f$
vanish linearly.  At $U\to1$, $p=O(1-U)$ while $f$ grows only as
$\ln(1-U)$, and $d\to1-\beta^2>0$.  There is therefore no boundary
constant to supply from a different calculation.

Use the previous rational variable
$U=(w-w^{-1})/(2\sqrt X)$.  The remaining root factor cancels from
its differential,
\begin{equation}
 \frac{\dd U}{k}=\sqrt{\frac sX}\frac{\dd w}{w}.
\end{equation}
Writing $z=w-1$, define four rational kernels
\begin{align}
 \widetilde R_0&=\frac4{\beta E^2}\frac{U'\overline h_0}{d},&
 \widetilde R_1&=\frac4{\beta E^2}\frac{U'h_1}{d},\\
 \widetilde R_k&=\frac4{\beta E^2}\sqrt{\frac sX}\frac1w
 \left[\frac{b_1}d+\frac{16sU^2}{X(X+2)d^2}\right],&
 \widetilde R_f&=-\frac8{\beta E^2}\sqrt{\frac sX}
 \frac{U(1-U^2)}{wd^2}.
\end{align}
Here $U'=\dd U/\dd z$ and
\begin{equation}
 h_1=-\frac{X^3-6X^2-16X-8}{X^2}
 +\frac{(1-k^2)(X^2+2X+4)}{2X}
 -\frac{3X}{8}(1-k^2)^2.
\end{equation}
The logarithms are $L_\beta=\operatorname{artanh}(\beta U)$,
$L_1=\operatorname{artanh}U$, and $L_k=\operatorname{artanh}(U/k)$.
The transformed integrand is
\begin{equation}
 \widetilde R_0 L_\beta+\widetilde R_1L_\beta L_1
 +\widetilde R_kL_\beta L_k+\widetilde R_f L_k.
\end{equation}
Every denominator now factors into powers of $w$ and
\begin{equation}
 Q_w=-Xw^4+(4X^2+18X+16)w^2-X.
\end{equation}
The roots of $Q_w$ are the four real poles
$\pm r,\pm r^{-1}$, where
$r=\sqrt{1+X/\beta^2}-\sqrt X/\beta$.
The factor $1+w^2$ is absent.
Thus the entire GPL expression, including the letters introduced by
partial fractions, has only real letters.
